\chapter{2007年北京卷（文）}
\section{单选题}
\begin{problem}
已知 \(\cos \theta \cdot \tan \theta < 0\)，那么角 \(\theta\) 是（\quad）

\choices
    {第一或第二象限角}
    {第二或第三象限角}
    {第三或第四象限角}
    {第一或第四象限角}
\end{problem}

\begin{problem}
函数 \( f(x) = 3^{x} (0 < x \leqslant 2) \) 的反函数的定义域为（\quad）

\choices
    {\((0, +\infty)\)}
    {(1, 9]}
    {\((0,1)\)}
    {\([9, +\infty)\)}
\end{problem}

\begin{problem}
函数 \( f(x) = \sin 2x - \cos 2x \) 的最小正周期是（\quad）

\choices
    {\(\frac{\pi}{2}\)}
    {\(\pi\)}
    {\(2\pi\)}
    {\(4\pi\)}
\end{problem}

\begin{problem}
椭圆 \(\frac{x^2}{a^2} + \frac{y^2}{b^2} = 1 (a > b > 0)\) 的焦点为 \(F_1, F_2\)，两条准线与 \(x\) 轴的交点分别为 \(M, N\). 若 \(|MN| \leqslant 2|F_1F_2|\)，则该椭圆离心率的取值范围是（\quad）

\choices
    {\(\left(0, \frac{1}{2}\right]\)}
    {\(\left(0, \frac{\sqrt{2}}{2}\right]\)}
    {\(\left[\frac{1}{2}, 1\right)\)}
    {\(\left[\frac{\sqrt{2}}{2}, 1\right)\)}
\end{problem}

\begin{problem}
某城市的汽车牌照号码由 2 个英文字母后接 4 个数字组成. 其中 4 个数字互不相同的牌照号码共有（\quad）

\choices
    {\((\mathrm{C}_{50}^{1})^{4}\mathrm{A}_{50}^{1}\) 个}
    {\(\mathrm{A}_{50}^{2}\mathrm{A}_{50}^{1}\) 个}
    {\((\mathrm{C}_{50}^{1})^{4}10^{4}\) 个}
    {\(\mathrm{A}_{50}^{3}10^{4}\) 个}
\end{problem}

\begin{problem}
若不等式组 \( \left\{\begin{aligned}x-y+5&\geqslant0,\\ y&\geqslant a,\\ 0&\leqslant x\leqslant2\end{aligned}\right. \) 表示的平面区域是一个三角形，则 a 的取值范围是（\quad）

\choices
    {\(a < 5\)}
    {\(a \geqslant 7\)}
    {\(5 \leqslant a < 7\)}
    {\(a < 5\) 或 \(a \geqslant 7\)}
\end{problem}

\begin{problem}
平面 \(\alpha //\) 平面 \(\beta\) 的一个充分条件是（\quad）

\choices
    {存在一条直线 \(a, a // \alpha, a // \beta\)}
    {存在一条直线 \(a, a \subset \alpha, a // \beta\)}
    {存在两条平行直线 \(a, b, a \subset \alpha, b < \beta, a // \beta, b // \alpha\)}
    {存在两条异面直线 \(a, b, a \subset \alpha, b < \beta, a // \beta, b // \alpha\)}
\end{problem}

\begin{problem}
对于函数\circled{1} \( f(x) = |x + 2| \)，\circled{2} \( f(x) = (x - 2)^2 \)，\circled{3} \( f(x) = \cos (x - 2) \)，判断如下两个命题的真假： 命题甲： \( f(x+2) \) 是偶函数； 命题乙： \( f(x) \) 在 \( (-\infty,2) \) 上是减函数，在 \( (2,+\infty) \) 上是增函数； 能使命题甲、乙均为真的所有函数的序号是（\quad）

\choices
    {\circled{1}\circled{2}}
    {\circled{1}\circled{3}}
    {\circled{2}}
    {\circled{3}}
\end{problem}

\section{填空题}
\begin{problem}
\(f(x)\) 是 \(f(x) = \frac{1}{3} x^3 + 2x + 1\) 的导函数，则 \(f'(-1)\) 的值是 \fillinblank{} \fillinblank{}.
\end{problem}

\begin{problem}
若数列 \(\{a_{n}\}\) 的前 \(n\) 项和 \(S_{n} = n^{2} - 10n (n = 1,2,3,\dots)\)，则此数列的通项公式为 \fillinblank{} \fillinblank{}.
\end{problem}

\begin{problem}
已知向量 \(a = (2,4)\)，\(b = (1,1)\). 若向量 \(b \perp (a + \lambda b)\)，则实数 \(\lambda\) 的值是 \fillinblank{} \fillinblank{}.
\end{problem}

\begin{problem}
在 \(\triangle ABC\) 中，若 \(\tan A = \frac{1}{3}, C = 150^{\circ}, BC = 1\)，则 \(AB = \fillinblank{}\).
\end{problem}

\begin{problem}
2002 年在北京召开的国际数学家大会，会标是以我国古代数学家赵爽的弦图为基础设计的. 弦图是由四个全等直角三角形与一个小正方形拼成的一个大正方形 (如图). 如果小正方形的面积为 1，大正方形的面积为 25，直角三角形中较小的锐角为 \(\theta\)，那么 \(\cos 2\theta\) 的恒等于 \fillinblank{} \fillinblank{}.
\end{problem}

\begin{problem}
已知函数 \( f(x), g(x) \) 分别由下表给出 x 1 2 3 f(x) 2 1 1 则 \( f[g(1)] \) 的值为 \fillinblank{}；当 \( g[f(x)] = 2 \) 时，z =\fillinblank{} \fillinblank{}.
\end{problem}

\section{解答题}
\begin{problem}
定义于 \(x\) 的不等式 \(\frac{x - a}{x + 1} < 0\) 的解集为 \(P\)，不等式 \(|x - 1| \leqslant 1\) 的解集为 \(Q\). (1) 若 \(a = 3\)，求 \(P\); (2) 若 \(Q \subseteq P\)，求正数 \(a\) 的取值范围.
\end{problem}

\begin{problem}
数列 \(\{a_{n}\}\) 中，\(a_{1} = 2\)，\(a_{n + 1} = a_{n} + cn\) ( \(c\) 是常数，\(n = 1,2,3,\dots\))，且 \(a_{1}\)，\(a_{2}, a_{3}\) 成公比不为 1 的等比数列. (1) 求 \(c\) 的值; (2) 求 \(\{a_{n}\}\) 的通项公式.
\end{problem}

\begin{problem}
如图，在 Rt\(\triangle\)AOB 中，\(\angle\)OAB = \( \frac{\pi}{6} \) ，斜边 AB = 4. Rt\(\triangle\)AOC 可以通过 Rt\(\triangle\)AOB 以直线 AO 为轴旋转得斜，且二面角 B - AO - C 是直二面角. D 是 AB 的中点. (1) 求证: 平面 \(COD \perp\) 平面 \(AOB\); (2) 求异面直线 AO 与 CD 所成角的大小.
\end{problem}

\begin{problem}
某条公共汽车线路沿线共有 11 个车站（包括起点站和终点站）. 在起点站开出的一辆公共汽车上有 6 位乘客，假设每位乘客在起点站之外的各个车站下车是等可能的，求： (1) 这 6 位乘客在互不相同的车站下车的概率; (2) 这 6 位乘客中恰有 3 人在终点站下车的概率;
\end{problem}

\begin{problem}
如图，矩形 \(ABCD\) 的两条对角线相交于点 \(M(2,0)\)，\(AB\) 边所在直线的方程为 \(x - 3y - 6 = 0\)，点 \(\mathcal{T}(-1,1)\) 在 \(AD\) 边所在直线上. (1) 求 AD 边所在直线的方程; (2) 求矩形 ABCD 外接圆的方程; (3) 若动圆 P 过点 \( N(-2,0) \) ，且与矩形 ABCD 的外接圆外切，求动圆 P 的圆心的轨迹方程.
\end{problem}

\begin{problem}
已知函数 \( y = kx \) 与 \( y = x^2 + 2 (x \geqslant 0) \) 的图象相交于 \( A(x_1, y_1), B(x_2, y_2) \). \( l_1, l_2 \) 分别是 \( y = x^2 + 2 (x \geqslant 0) \) 的图象在 \( A, B \) 两点的切线，\( M, N \) 分别是 \( l_1, l_2 \) 与 \( x \) 轴的交点. (1) 求 k 的取值范围; (2) 设 \(t\) 为点 \(M\) 的横坐标，当 \(x_{1} < x_{2}\) 时，写出 \(t\) 以 \(x_{1}\) 为自变量的函数式，并求其定义域和值域; (3) 试比较 \( \{OM\} \) 与 \( \{ON\} \) 的大小，并说明理由 (O 是坐标原点).
\end{problem}

